\documentclass[11pt]{article}
\usepackage[a4paper,margin=1in]{geometry}
\usepackage{fontspec}
\usepackage[boldfont=false]{xeCJK}
\setCJKmainfont{FandolSong-Regular.otf}
\usepackage{amsmath}
\usepackage{amssymb}
\usepackage{array}
\usepackage{booktabs}
\usepackage{graphicx}
\usepackage{adjustbox}
\usepackage{enumitem}
\setlist[itemize]{topsep=2pt,partopsep=0pt,parsep=0pt,itemsep=2pt,leftmargin=1.4em}
\setlist[enumerate]{topsep=2pt,partopsep=0pt,parsep=0pt,itemsep=2pt,leftmargin=1.6em}
\usepackage{parskip}
\usepackage{fvextra}
\fvset{breaklines=true,breakanywhere=true,fontsize=\small}
\setlength{\parindent}{0pt}

\begin{document}

\begin{center}
  {\LARGE\bfseries Acids, bases and salts}\\[0.35em]
  {\large Igcse Chemistry}
\end{center}
\vspace{0.6em}

\section*{Acids}

An \textbf{acid} 酸 is a substance that forms hydrogen ions ($\text{H}^{+}$) when dissolved in water. Acids have three typical reactions:

\begin{itemize}
\item with \textbf{metals} 金属: acid + metal → a \textbf{salt} 盐 + \textbf{hydrogen} 氢气

\item with bases: acid + \textbf{base} 碱 → a salt + water (this is \textbf{neutralisation} 中和)

\item with \textbf{carbonates} 碳酸盐: acid + carbonate → a salt + water + \textbf{carbon dioxide} 二氧化碳

\end{itemize}

\subsection*{Indicators}

An \textbf{indicator} 指示剂 is a dye that changes colour to show whether a solution is acidic or alkaline.

\begin{center}
\adjustbox{max width=\linewidth}{%
\begin{tabular}{lll}
\toprule
\textbf{Indicator} & \textbf{In acid} & \textbf{In alkali} \\
\midrule
\textbf{litmus} 石蕊 & red & blue \\
\textbf{thymolphthalein} 百里酚酞 & colourless & blue \\
\textbf{methyl orange} 甲基橙 & red & yellow \\
\bottomrule
\end{tabular}}
\end{center}

Aqueous solutions of acids contain hydrogen \textbf{ions} 离子 ($\text{H}^{+}$). Aqueous solutions of alkalis contain hydroxide ions ($\text{OH}^{-}$).

\section*{Bases and alkalis}

Bases are \textbf{oxides} 氧化物 or \textbf{hydroxides} 氢氧化物 of metals. An \textbf{alkali} 可溶性碱 is a base that dissolves in water (a soluble base).

Bases have two typical reactions:

\begin{itemize}
\item with acids: base + acid → a salt + water (neutralisation again)

\item with \textbf{ammonium salts} 铵盐: this releases ammonia gas.

\end{itemize}

\section*{Strong and weak acids}

An acid is a \textbf{proton} 质子 donor — it gives away $\text{H}^{+}$ ions (a hydrogen ion is just a proton). A base is a proton acceptor.

How strong an acid is depends on how much of it splits into ions in water:

\begin{itemize}
\item A \textbf{strong acid} 强酸 is completely \textbf{dissociated} 电离 (fully split into ions) in water. \textbf{Hydrochloric acid} 盐酸 is strong:

\end{itemize}

\[
\text{HCl}(aq) \rightarrow \text{H}^{+}(aq) + \text{Cl}^{-}(aq)
\]

\begin{itemize}
\item A \textbf{weak acid} 弱酸 is only partly dissociated. \textbf{Ethanoic acid} 乙酸 is weak, so the equation uses the reversible arrow:

\end{itemize}

\[
\text{CH}_3\text{COOH}(aq) \rightleftharpoons \text{H}^{+}(aq) + \text{CH}_3\text{COO}^{-}(aq)
\]

Note: 'strong' and 'weak' are about dissociation, not about being dilute or concentrated.

\section*{The pH scale}

The pH scale runs from 0 to 14 and tells you how acidic or alkaline a solution is. A higher hydrogen ion concentration means a lower pH.

You can find pH using \textbf{universal indicator} 通用指示剂, which turns different colours:

\begin{center}
\adjustbox{max width=\linewidth}{%
\begin{tabular}{lll}
\toprule
\textbf{pH} & \textbf{Type} & \textbf{Colour of universal indicator} \\
\midrule
below 7 & \textbf{acidic} 酸性 & red / orange / yellow \\
exactly 7 & \textbf{neutral} 中性 & green \\
above 7 & \textbf{alkaline} 碱性 & blue / purple \\
\bottomrule
\end{tabular}}
\end{center}

When an acid and an alkali react, the $\text{H}^{+}$ and $\text{OH}^{-}$ ions join to make water:

\[
\text{H}^{+}(aq) + \text{OH}^{-}(aq) \rightarrow \text{H}_2\text{O}(l)
\]

\section*{Oxides}

Oxides can be sorted by how they behave:

\begin{itemize}
\item \textbf{Acidic oxides} are oxides of non-metals, such as $\text{SO}_2$ and $\text{CO}_2$.

\item \textbf{Basic oxides} are oxides of metals, such as $\text{CuO}$ and $\text{CaO}$.

\item \textbf{Amphoteric} 两性 oxides react with both acids and bases to make a salt and water. $\text{Al}_2\text{O}_3$ and $\text{ZnO}$ are amphoteric.

\end{itemize}

So metal oxides tend to be basic and non-metal oxides tend to be acidic.

\section*{Preparing salts}

Whether a salt can be made by a certain method depends on whether it is \textbf{soluble} 可溶 (dissolves) or \textbf{insoluble} 不溶 (does not dissolve).

\subsection*{Solubility rules}

\begin{center}
\adjustbox{max width=\linewidth}{%
\begin{tabular}{ll}
\toprule
\textbf{Salt type} & \textbf{Rule} \\
\midrule
sodium, potassium, ammonium salts & all soluble \\
\textbf{nitrates} 硝酸盐 & all soluble \\
\textbf{chlorides} 氯化物 & soluble, except lead and silver \\
\textbf{sulfates} 硫酸盐 & soluble, except barium, calcium and lead \\
carbonates & insoluble, except sodium, potassium and ammonium \\
hydroxides & insoluble, except sodium, potassium, ammonium and (partly) calcium \\
\bottomrule
\end{tabular}}
\end{center}

\subsection*{Making a soluble salt}

You react an acid with one of these:

\begin{itemize}
\item an \textbf{alkali}, using \textbf{titration} 滴定 (since both are solutions, you must measure the exact volumes);

\item an \textbf{excess} 过量 of a metal, an insoluble base, or an insoluble carbonate.

\end{itemize}

When you use an excess of a solid, you then \textbf{filter} 过滤 to remove the leftover solid. Finally you \textbf{evaporate} 蒸发 some water and let the solution \textbf{crystallise} 结晶 to get the salt.

\subsection*{Making an insoluble salt}

An insoluble salt is made by \textbf{precipitation} 沉淀: mix two solutions that each contain one of the needed ions, and the insoluble salt forms as a solid. You then filter, wash and dry it.

\subsection*{Water in salts}

\begin{itemize}
\item A \textbf{hydrated} 水合 substance is chemically joined with water.

\item An \textbf{anhydrous} 无水 substance contains no water.

\end{itemize}

The water molecules inside hydrated crystals are called the \textbf{water of crystallisation} 结晶水. For example, $\text{CuSO}_4 \bullet 5\text{H}_2\text{O}$ has five water molecules for each formula unit.


\end{document}
